\begin{papercard}{hassan2018}
\cf{Stages} \stagetag{H} (a concrete pathway: public activity data plus a protection feature that leaks the protected location), \stagetag{C} (a fitness-specific privacy feature whose claimed protection can be tested against a known attack, for Axis C). USENIX Security 2018 paper (University of Illinois). Added from the gap search of 23 October 2026.

\cf{Setting} Fitness social networks publish routes that combine location and health data; besides general settings they offer endpoint privacy zones (EPZ) that hide activity near a sensitive location \pc{hassan2018}{2}{domain-specific mechanisms such as Endpoint Privacy Zones (EPZs) conceal fitness activity that occurs within a certain distance of sensitive user locations such as homes or work places}. Whether such features are used, and whether they work, was unknown \pc{hassan2018}{2}{nor is it clear that these countermeasures are adequate to prevent attackers from discovering the private locations of users}.

\cf{Data and findings} A month of public Strava posts (21 million posts, 3 million users): 16.5\,\% of users had an EPZ \pc{hassan2018}{2}{we observe that 16.5\% of users make use of Endpoint Privacy Zones (EPZs)}. An attack using the remaining route information located the protected place for most active users \pc{hassan2018}{2}{discovering that 95.1\% of moderately active users are at risk of having their protected locations ex- tracted by an attacker}. The attack generalised to Garmin Connect and Map My Tracks.

\cf{Countermeasures and disclosure} The authors tested geo-indistinguishability and proposed EPZ fuzzing; the vendors acknowledged the problem and adopted countermeasures \pc{hassan2018}{3}{All companies have acknowledged the vulnerability and have incor- porated one or more of our proposed countermeasures into their production systems}.

\cf{Limits (card)} 2016--2017 platform behaviour, since changed; the attack concerns public posts; the paper does not analyse how vendors describe the feature to users. Dhondt et al.\ \pp{dhondt2022} show the fixes were incomplete.

\cf{Use for this thesis} (1) Harm tree (fitness): a documented pathway from route data through an adversary with public access to a physical-location consequence (stalking, burglary), with the protection's failure condition named. (2) Axis C test case: if a sampled app (Strava, Garmin Connect, adidas Running, Nike Run Club are in the Play pilot list) describes a privacy zone as hiding the home location, the claim's scope can be compared with this pathway; this is the kind of ``potential mismatch'' the methodology asks a human to judge. (3) Shows that adoption of protective features is low (one in six), which matters for responsibility-shifting stance (Axis B).
\end{papercard}
